\subsection{Pregled naučnoistraživačkih radova}\label{sec:radovi}
Svaki od pet članova analizira najmanje jedan relevantan naučnoistraživački rad.
Za svaki rad navode se \textbf{problem, skup podataka, metode, metrike i najbolji rezultat}.
Tri rada (\cite{afriyie2023fraud,albulushi2026fraud,breskuviene2024fraud}) koriste isti
Sparkov skup podataka kao naš projekat, a dva (\cite{ileberi2022fraud,sulaiman2024fraud})
najčešće korišteni stvarni evropski skup transakcija, što omogućava poređenje nalaza na
simuliranim i stvarnim podacima.


\Needspace{12\baselineskip}
\subsubsection*{Naučnoistraživački rad -- \ClanJedan}

\textbf{Naslov rada:}
\emph{A supervised machine learning algorithm for detecting
and predicting fraud in credit card transactions}.\\

\textbf{Autori:}
Jonathan Kwaku Afriyie, Kassim Tawiah, Wilhemina Adoma Pels,
Sandra Addai-Henne, Harriet Achiaa Dwamena, Emmanuel Odame Owiredu,
Samuel Amening Ayeh i John Eshun.\\

\textbf{Izvor:}
\emph{Decision Analytics Journal}, vol. 6, članak 100163,
2023. \cite{afriyie2023fraud}

\begin{table}[H]
\centering
\caption{Sažeti pregled naučnoistraživačkog rada Afriyie i saradnika (2023)}
\label{tab:afriyie-pregled}
\small
\renewcommand{\arraystretch}{1.22}
\begin{tabularx}{\linewidth}{@{}p{0.25\linewidth}Y@{}}
\toprule
\textbf{Stavka} & \textbf{Opis} \\
\midrule

Problem &
Detekcija prevarnih transakcija kreditnim karticama kao problem
binarne klasifikacije, pri čemu se razlikuju legitimne i
prevarne transakcije. \\

Skup podataka &
\emph{Credit Card Transactions Fraud Detection Dataset}
(Kaggle), simuliran pomoću alata Sparkov.
U radu se navodi 555\,719 transakcija i 23 kolone
u početnom skupu podataka. \\

Korištene metode &
Logistička regresija (Logistic Regression),
stablo odlučivanja (Decision Tree) i slučajna šuma
(Random Forest). Priprema podataka uključuje
skaliranje i nasumično smanjivanje većinske klase
(undersampling). \\

Metrike &
Tačnost (accuracy), preciznost (precision),
osjetljivost (recall/sensitivity), specifičnost,
F1-mjera i ROC-AUC. \\

Najbolji rezultat &
\textbf{Random Forest}: tačnost \textbf{96\%},
ROC-AUC \textbf{98,9\%}.
Prijavljene vrijednosti: recall 0,97,
precision 0,09 i F1-mjera 0,17. \\

\bottomrule
\end{tabularx}
\end{table}

\textbf{Cilj i istraživački problem.}

Autori istražuju koliko uspješno različiti algoritmi nadziranog
mašinskog učenja razlikuju prevarne od legitimnih transakcija.
Osnovni cilj jeste uporediti performanse tri klasifikacijska
algoritma i utvrditi koji je najpogodniji za analizirani
skup podataka.

Dodatno istražuju obrasce povezane s iznosom i vremenom
transakcije, kategorijom trgovca te pojedinim
karakteristikama korisnika.

\textbf{Skup podataka i njegova nebalansiranost.}

Rad koristi javno dostupan Kaggle skup simuliranih transakcija,
obavljenih tokom 2020. godine u Sjedinjenim Američkim Državama.

Skup obuhvata numeričke i kategoričke atribute,
a ciljna varijabla predstavlja status prevare.

Nakon postupaka obrade, u Tabeli 4 rada prikazano je
482\,672 legitimnih i 2\,135 prevarnih transakcija,
odnosno približno 99,6\% i 0,4\%.

Takva raspodjela predstavlja izraženu nebalansiranost klasa,
pa sama ukupna tačnost nije dovoljna za procjenu
uspješnosti klasifikatora.

\textbf{Priprema podataka i korištene metode.}

Autori navode čišćenje podataka, transformaciju numeričkih
atributa na raspon od 0 do 1 i primjenu metode
\emph{random undersampling}.

Ova metoda smanjuje broj primjera većinske klase kako bi
se ublažila pristrasnost modela prema legitimnim transakcijama.

Na Slici 1 rada prikazan je uravnotežen uzorak s
1\,708 prevarnih i 1\,707 legitimnih transakcija.

Za klasifikaciju se zatim koriste logistička regresija,
jedno stablo odlučivanja i Random Forest,
koji kombinuje veći broj stabala.

\textbf{Poređenje eksperimentalnih rezultata.}

U Tabeli~\ref{tab:afriyie-rezultati} prikazane su vrijednosti
koje su autori objavili u Tabeli 11 i uz ROC krive (Slika 16).

\begin{table}[H]
\centering
\caption{Poređenje rezultata klasifikacijskih algoritama iz rada}
\label{tab:afriyie-rezultati}
\small
\setlength{\tabcolsep}{4pt}
\renewcommand{\arraystretch}{1.18}

\begin{tabularx}{\linewidth}{@{}Y c c c c c c@{}}
\toprule

\textbf{Model} &
\textbf{Acc.} &
\textbf{Prec.} &
\textbf{Recall} &
\textbf{F1} &
\textbf{Spec.} &
\textbf{AUC} \\

\midrule

Decision Tree &
0,92 & 0,05 & 0,93 & 0,09 & 0,92 & 0,945 \\

Logistic Regression &
0,92 & 0,04 & 0,76 & 0,08 & 0,92 & 0,872 \\

\textbf{Random Forest} &
\textbf{0,96} &
\textbf{0,09} &
\textbf{0,97} &
\textbf{0,17} &
\textbf{0,96} &
\textbf{0,989} \\

\bottomrule
\end{tabularx}
\end{table}

\textbf{Tumačenje rezultata.}

Random Forest je postigao najveće vrijednosti svih
prikazanih metrika.

Njegov ROC-AUC od 0,989 ukazuje na veoma dobru sposobnost
razlikovanja klasa pri promjeni klasifikacijskog praga.

Međutim, precision od samo 0,09 znači da je tek približno
9\% transakcija označenih kao prevarne zaista bilo prevarno.

Iz objavljene konfuzijske matrice (Tabela 7 rada)
vidi se da je model prepoznao 409 prevara i propustio 18,
ali je 4\,052 legitimne transakcije pogrešno označio
kao prevarne.

Zato visoka tačnost i ROC-AUC ne znače da model proizvodi
malo lažnih alarma.

Niska F1-mjera (0,17) dodatno ukazuje na neravnotežu
između preciznosti i osjetljivosti modela.

\textbf{Dodatna zapažanja iz istraživačke analize.}

Autori izvještavaju da se veći broj prevarnih transakcija
u njihovom skupu javlja između 22:00 i 04:00 GMT te
posebno razmatraju starosnu grupu korisnika iznad
60 godina i kategoriju kupovine.

Budući da su podaci simulirani, ova zapažanja
predstavljaju obrasce uočene u posmatranom skupu,
a ne dokazane zakonitosti stvarnog finansijskog poslovanja.

\textbf{Zaključak i značaj za istraživanje.}

Istraživanje upućuje na to da Random Forest može biti
uspješan izbor za detekciju prevarnih transakcija,
posebno prema ROC-AUC i osjetljivosti.

Istovremeno pokazuje da izrazita nebalansiranost klasa
zahtijeva pažljivu evaluaciju i praćenje lažnih alarma.

Nalazi rada predstavljaju osnovu za očekivanje da će
metode zasnovane na stablima ostvarivati konkurentne rezultate,
ali i da će precision, recall i F1-mjera biti važnije
za praktičnu interpretaciju nego sama ukupna tačnost.


\Needspace{12\baselineskip}
\subsubsection*{Naučnoistraživački rad -- \ClanDva}

\textbf{Naslov rada:}
\emph{Applying supervised machine learning algorithms and ensemble
models to enhance credit card fraud detection}.\\

\textbf{Autori:}
Abrar Al-Bulushi, Abdul K. Shaikh i Naresh Adhikari.\\

\textbf{Izvor:}
\emph{Frontiers in Artificial Intelligence}, vol. 9, članak 1813728,
2026. \cite{albulushi2026fraud}

\begin{table}[H]
\centering
\caption{Sažeti pregled naučnoistraživačkog rada Al-Bulushi i saradnika (2026)}
\label{tab:albulushi-pregled}
\small
\renewcommand{\arraystretch}{1.22}
\begin{tabularx}{\linewidth}{@{}p{0.25\linewidth}Y@{}}
\toprule
\textbf{Stavka} & \textbf{Opis} \\
\midrule

Problem &
Detekcija prevarnih transakcija kreditnim karticama kao binarna
klasifikacija na izrazito nebalansiranim podacima, uz provjeru
generalizacije modela na neviđenim skupovima podataka. \\

Skup podataka &
Glavni skup za treniranje i testiranje je \emph{isti Sparkov skup}
koji koristi naš tim (približno 1,3 miliona transakcija, 0,5\,\% prevara,
22 atributa). Skup od 555\,719 transakcija iz iste kolekcije i još četiri
javna skupa (među njima evropski skup od 284\,807 transakcija i PaySim)
služe kao neviđeni skupovi. \\

Korištene metode &
Logistička regresija, Naive Bayes, stablo odlučivanja, Random Forest,
neuronska mreža (ANN) i XGBoost; uravnotežavanje metodama RUS, ROS
i SMOTE; selekcija atributa (korelacija, ANOVA, RFE, unaprijedna
selekcija); novi atributi ponašanja korisnika; pomjeranje praga
odlučivanja; ansambli \emph{bagging}, \emph{boosting} i \emph{stacking}. \\

Metrike &
Tačnost, preciznost, osjetljivost (recall), F1-mjera i AUC,
uz konfuzijsku matricu; 4-struka unakrsna validacija. \\

Najbolji rezultat &
Pojedinačno: \textbf{XGBoost} bez uravnotežavanja, F1 \textbf{0,84}
(AUC 0,99). Konačni model \textbf{Bagging~1} (10 modela):
F1 \textbf{0,795}, recall 0,864, AUC 0,99; uz nove atribute ponašanja
F1 raste na \textbf{0,83} (recall 0,89). \\

\bottomrule
\end{tabularx}
\end{table}

\textbf{Cilj i istraživački problem.}
Autori razvijaju cjelovit okvir za detekciju prevara koji objedinjuje
izbor algoritma, uravnotežavanje klasa, selekciju i konstrukciju atributa
te ansambl modele. Posebno ih zanima da li model treniran na jednom
skupu zadržava performanse na drugim, neviđenim skupovima podataka.

\textbf{Skup podataka i priprema.}
Rad je za nas posebno relevantan jer kao glavni skup koristi upravo
Sparkov podatke sa Kaggle platforme, sa 12 kategoričkih i 11 numeričkih
atributa i bez nedostajućih vrijednosti. Podaci su standardizovani
(\emph{z-score}), a podjela je urađena 4-strukom unakrsnom validacijom
(75\,\% podataka za treniranje). Autori su konstruisali i nove atribute
ponašanja: oznaku da li je transakcija izvan uobičajenog vremenskog prozora
korisnika (izvan 10.--90.~percentila sati transakcija) i \emph{anomaly score}
dobijen metodom Isolation Forest.

\textbf{Poređenje eksperimentalnih rezultata.}
Tabela~\ref{tab:albulushi-rezultati} prikazuje izdvojene rezultate iz
Tabela 1, 8 i 9 rada (pojedinačni modeli, pomjeranje praga i ansambli).

\begin{table}[H]
\centering
\caption{Izdvojeni rezultati iz rada Al-Bulushi i saradnika (Sparkov skup)}
\label{tab:albulushi-rezultati}
\small
\setlength{\tabcolsep}{4pt}
\renewcommand{\arraystretch}{1.18}
\begin{tabularx}{\linewidth}{@{}Y c c c c c@{}}
\toprule
\textbf{Model} & \textbf{Acc.} & \textbf{F1} & \textbf{Prec.} & \textbf{Recall} & \textbf{AUC} \\
\midrule
Logistička regresija -- RUS & 0,87 & 0,06 & 0,03 & 0,76 & 0,96 \\
Naive Bayes -- bez uravnoteženja & 0,90 & 0,07 & 0,04 & 0,68 & 0,82 \\
Stablo odlučivanja -- ROS & 0,99 & 0,78 & 0,79 & 0,77 & 0,88 \\
Random Forest -- bez uravnoteženja & 0,99 & 0,82 & 0,88 & 0,76 & 0,97 \\
Random Forest -- RUS & 0,97 & 0,30 & 0,17 & 0,97 & 0,99 \\
ANN -- bez uravnoteženja & 0,99 & 0,83 & 0,93 & 0,76 & 0,99 \\
\textbf{XGBoost -- bez uravnoteženja} & \textbf{0,99} & \textbf{0,84} & \textbf{0,91} & \textbf{0,79} & \textbf{0,99} \\
XGBoost -- SMOTE, prag 0,7 & 0,99 & 0,70 & 0,67 & 0,85 & 0,99 \\
\midrule
Bagging~1 (10 modela) & 0,99 & 0,79 & 0,73 & 0,86 & 0,99 \\
Stacking~1 (RF baza, ANN meta) & 0,99 & 0,80 & 0,79 & 0,80 & 0,93 \\
Boosting~2 (LR + AdaBoost) & 0,99 & 0,00 & 0,00 & 0,00 & 0,82 \\
\bottomrule
\end{tabularx}
\end{table}

\textbf{Tumačenje rezultata.}
Logistička regresija i Naive Bayes daju vrlo nisku F1-mjeru (ispod 0,1)
bez obzira na metodu uravnotežavanja, dok modeli zasnovani na stablima
(Random Forest, XGBoost) i neuronska mreža postižu F1 oko 0,8.
Rezultati jasno pokazuju kompromis između preciznosti i osjetljivosti:
\emph{random undersampling} podiže recall Random Foresta na 0,97, ali
preciznost pada na 0,17, jer model treniran na uravnoteženim podacima
na realističnoj raspodjeli generiše mnogo lažnih alarma. Tačnost je
za gotovo sve modele 0,99 i zato nije informativna. Ansambl Bagging~1
nije postigao najveći F1, ali su ga autori izabrali kao konačni model zbog
stabilnosti i visokog recall-a, a na neviđenom Sparkov skupu postiže
F1 0,61, odnosno 0,66 uz nove atribute ponašanja.

\textbf{Zaključak i značaj za istraživanje.}
Rad potvrđuje da su na našem skupu podataka metode zasnovane na
stablima i \emph{boosting} najperspektivnije, a da su linearni modeli i
Naive Bayes slabi. Za naš projekat su posebno korisna tri zapažanja:
(1) uravnotežavanje treba primjenjivati samo na skupu za treniranje i
pratiti njegov uticaj na preciznost, (2) pomjeranje praga odlučivanja je
jednostavan način da se podesi odnos preciznosti i osjetljivosti, i
(3) atributi izvedeni iz vremena transakcije i ponašanja korisnika
mogu dodatno poboljšati detekciju.


\Needspace{12\baselineskip}
\subsubsection*{Naučnoistraživački rad -- \ClanTri}

\textbf{Naslov rada:}
\emph{Enhancing credit card fraud detection: highly imbalanced data case}.\\

\textbf{Autori:}
Dalia Breskuvienė i Gintautas Dzemyda.\\

\textbf{Izvor:}
\emph{Journal of Big Data}, vol. 11, članak 182,
2024. \cite{breskuviene2024fraud}

\begin{table}[H]
\centering
\caption{Sažeti pregled naučnoistraživačkog rada Breskuvienė i Dzemyde (2024)}
\label{tab:breskuviene-pregled}
\small
\renewcommand{\arraystretch}{1.22}
\begin{tabularx}{\linewidth}{@{}p{0.25\linewidth}Y@{}}
\toprule
\textbf{Stavka} & \textbf{Opis} \\
\midrule

Problem &
Selekcija atributa za detekciju prevara kreditnim karticama kada su
podaci izrazito nebalansirani; uobičajene metode selekcije atributa
zanemaruju manjinsku klasu. \\

Skup podataka &
Tri skupa: sintetički skup transakcija (3\,445\,553 instanci,
0,14\,\% prevara), \emph{Sparkov} skup -- spojene datoteke
\texttt{fraudTrain} i \texttt{fraudTest} (1\,852\,394 instanci,
0,52\,\% prevara) i stvarni evropski skup (284\,807 instanci,
0,17\,\% prevara). \\

Korištene metode &
Nova metoda selekcije atributa FID-SOM zasnovana na samoorganizujućim
mapama (SOM), upoređena s univarijatnom selekcijom (F-test, $\chi^2$,
uzajamna informacija), RFE i važnošću atributa iz XGBoosta.
Klasifikatori: Random Forest, XGBoost i CatBoost. \\

Metrike &
F1-mjera, MCC, G-Mean, AUC-PR i AUC-ROC; svaki rezultat je prosjek
pet ponavljanja s različitim \emph{random seed} vrijednostima. \\

Najbolji rezultat &
Na Sparkov skupu: \textbf{XGBoost + FID-SOM} (7 atributa),
F1 \textbf{0,855}, AUC-PR 0,895, AUC-ROC 0,998, MCC 0,857
(bez selekcije atributa F1 je 0,543). \\

\bottomrule
\end{tabularx}
\end{table}

\textbf{Cilj i istraživački problem.}
Autori polaze od toga da su podaci o prevarama nebalansirani, visoko
dimenzionalni (posebno nakon kodiranja kategoričkih atributa) i podložni
promjeni obrazaca kroz vrijeme (\emph{concept drift}). Cilj rada je
pokazati da pažljivo odabran, manji skup atributa poboljšava klasifikaciju
više nego dodatni atributi.

\textbf{Skup podataka i priprema.}
Za nas je posebno važno što je drugi skup u radu isti Sparkov skup,
pri čemu su autori spojili datoteke za treniranje i testiranje i koristili
11 atributa. Umjesto nasumične podjele, podaci su podijeljeni
\emph{hronološki}: najranijih 80\,\% transakcija čini skup za treniranje, a
kasnijih 20\,\% skup za testiranje. Autori naglašavaju da nasumična
podjela vremenski zavisnih podataka daje previše optimistične rezultate.
Kategorički atributi su kodirani \emph{James-Stein} enkoderom, a numerički
normalizovani na raspon [0, 1]. Namjerno nije korišteno uravnotežavanje
klasa ni podešavanje hiperparametara, kako bi se izolovao čist efekat
selekcije atributa.

\textbf{Poređenje eksperimentalnih rezultata.}
Tabela~\ref{tab:breskuviene-rezultati} prikazuje rezultate iz Tabele 6 rada
za Sparkov skup, i to za model bez selekcije atributa (\emph{baseline})
i za predloženu metodu FID-SOM.

\begin{table}[H]
\centering
\caption{Rezultati na Sparkov skupu iz rada Breskuvienė i Dzemyde}
\label{tab:breskuviene-rezultati}
\small
\setlength{\tabcolsep}{4pt}
\renewcommand{\arraystretch}{1.18}
\begin{tabularx}{\linewidth}{@{}Y c c c c c@{}}
\toprule
\textbf{Model} & \textbf{F1} & \textbf{AUC-ROC} & \textbf{AUC-PR} & \textbf{MCC} & \textbf{G-Mean} \\
\midrule
CatBoost -- svi atributi & 0,750 & 0,991 & 0,769 & 0,765 & 0,791 \\
CatBoost -- FID-SOM & 0,828 & 0,997 & 0,866 & 0,831 & 0,870 \\
Random Forest -- svi atributi & 0,622 & 0,967 & 0,767 & 0,660 & 0,682 \\
Random Forest -- FID-SOM & 0,835 & 0,978 & 0,868 & 0,839 & 0,866 \\
XGBoost -- svi atributi & 0,543 & 0,973 & 0,641 & 0,593 & 0,622 \\
\textbf{XGBoost -- FID-SOM} & \textbf{0,855} & \textbf{0,998} & \textbf{0,895} & \textbf{0,857} & \textbf{0,889} \\
\bottomrule
\end{tabularx}
\end{table}

\textbf{Tumačenje rezultata.}
Na Sparkov skupu selekcija atributa ima veoma veliki efekat: F1-mjera
XGBoosta raste s 0,543 na 0,855, a Random Foresta s 0,622 na 0,835.
Istovremeno, AUC-ROC je visok (0,97--0,998) i za modele bez selekcije,
što pokazuje da AUC-ROC na izrazito nebalansiranim podacima slabo
razlikuje dobre i lošije modele. Zato autori stavljaju naglasak na
AUC-PR, MCC i F1-mjeru. Na stvarnom evropskom skupu, čiji su atributi
već dobijeni PCA transformacijom, efekat selekcije je mnogo manji
(F1 0,796 naspram 0,801 za XGBoost). Ukupno je FID-SOM bio najbolja
metoda u 32 od 45 poređenja (71\,\%).

\textbf{Zaključak i značaj za istraživanje.}
Rad pokazuje da su na našem skupu ansambli stabala (XGBoost, CatBoost,
Random Forest) uspješni, ali da njihove performanse snažno zavise od
izbora atributa. To za naš projekat znači da treba ukloniti identifikatore
i nebitne atribute, a kao metrike koristiti F1, AUC-PR i MCC, a ne
samo tačnost i AUC-ROC. Rad upozorava i da nasumična podjela podataka
može precijeniti performanse, pa je hronološka podjela vrijedna razmatranja.


\Needspace{12\baselineskip}
\subsubsection*{Naučnoistraživački rad -- \ClanCetiri}

\textbf{Naslov rada:}
\emph{A machine learning based credit card fraud detection using
the GA algorithm for feature selection}.\\

\textbf{Autori:}
Emmanuel Ileberi, Yanxia Sun i Zenghui Wang.\\

\textbf{Izvor:}
\emph{Journal of Big Data}, vol. 9, članak 24,
2022. \cite{ileberi2022fraud}

\begin{table}[H]
\centering
\caption{Sažeti pregled naučnoistraživačkog rada Ileberija i saradnika (2022)}
\label{tab:ileberi-pregled}
\small
\renewcommand{\arraystretch}{1.22}
\begin{tabularx}{\linewidth}{@{}p{0.25\linewidth}Y@{}}
\toprule
\textbf{Stavka} & \textbf{Opis} \\
\midrule

Problem &
Detekcija prevarnih transakcija kreditnim karticama uz izbor
optimalnog podskupa atributa pomoću genetskog algoritma (GA). \\

Skup podataka &
Evropski skup transakcija kreditnim karticama (Kaggle, septembar 2013.):
284\,807 transakcija, 0,172\,\% prevara, 30 numeričkih atributa
(V1--V28 dobijeni PCA transformacijom, \texttt{Time}, \texttt{Amount}).
Dodatna provjera na sintetičkom skupu s oko 24,4 miliona transakcija. \\

Korištene metode &
Genetski algoritam za selekciju atributa (funkcija prilagođenosti je
tačnost Random Foresta); klasifikatori: stablo odlučivanja, Random Forest,
logistička regresija, ANN i Naive Bayes; SMOTE za uravnotežavanje
i \emph{min-max} normalizacija. \\

Metrike &
Tačnost (glavna metrika), osjetljivost (recall), preciznost,
F1-mjera i AUC. \\

Najbolji rezultat &
\textbf{GA-RF} s vektorom atributa v5 (13 atributa): tačnost
\textbf{99,98\,\%}, preciznost 95,34\,\%, recall 72,56\,\%,
F1-mjera 82,41\,\%. \\

\bottomrule
\end{tabularx}
\end{table}

\textbf{Cilj i istraživački problem.}
Autori smatraju da veliki broj atributa može narušiti performanse
klasifikatora i predlažu sistem u kojem genetski algoritam bira
podskup atributa, a zatim se na tom podskupu trenira i upoređuje
pet klasičnih algoritama mašinskog učenja.

\textbf{Skup podataka i priprema.}
Korišten je najpoznatiji javni skup za ovaj problem: transakcije
evropskih korisnika kartica tokom dva dana, sa samo 492 prevare.
Za razliku od našeg skupa, svi atributi su numerički i anonimizovani.
Klasna nebalansiranost je ublažena metodom SMOTE, a atributi su
normalizovani \emph{min-max} postupkom. Genetski algoritam je
pronašao pet kandidatskih vektora atributa (v1--v5) dužine od 9 do
18 atributa, a atribut \texttt{Amount} se nalazi u svakom od njih.

\textbf{Poređenje eksperimentalnih rezultata.}
Tabela~\ref{tab:ileberi-rezultati} prikazuje rezultate za vektor v5
(Tabela 6 rada) i za vektor v1, na kojem je ostvarena najveća F1-mjera
(Tabela 2 rada).

\begin{table}[H]
\centering
\caption{Rezultati klasifikatora iz rada Ileberija i saradnika (u \%)}
\label{tab:ileberi-rezultati}
\small
\setlength{\tabcolsep}{4pt}
\renewcommand{\arraystretch}{1.18}
\begin{tabularx}{\linewidth}{@{}Y c c c c@{}}
\toprule
\textbf{Model (vektor atributa)} & \textbf{Acc.} & \textbf{Recall} & \textbf{Prec.} & \textbf{F1} \\
\midrule
\textbf{Random Forest (v5)} & \textbf{99,98} & \textbf{72,56} & \textbf{95,34} & \textbf{82,41} \\
Stablo odlučivanja (v5) & 99,89 & 72,56 & 65,07 & 68,61 \\
ANN (v5) & 99,08 & 77,87 & 12,27 & 21,20 \\
Naive Bayes (v5) & 99,44 & 57,52 & 15,85 & 24,85 \\
Logistička regresija (v5) & 99,77 & 46,90 & 34,64 & 39,84 \\
\midrule
Random Forest (v1) & 99,94 & 76,99 & 89,69 & 82,85 \\
ANN (v1) & 99,94 & 77,87 & 84,61 & 81,10 \\
Naive Bayes (v1) & 98,13 & 84,95 & 6,83 & 12,65 \\
\bottomrule
\end{tabularx}
\end{table}

\textbf{Tumačenje rezultata.}
Random Forest je najbolji ili među najboljima za svaki vektor atributa.
Iako svi modeli imaju tačnost iznad 98\,\%, preciznost i F1-mjera se
drastično razlikuju: Naive Bayes s vektorom v1 ima tačnost 98,13\,\%,
a preciznost samo 6,83\,\%, što znači da je velika većina transakcija
koje označi kao prevarne zapravo legitimna. Rezultati pokazuju i da
izbor atributa snažno utiče na pojedine modele (npr. F1 ANN-a pada
s 81,10\,\% na 21,20\,\% između vektora v1 i v5). Treba imati na umu
da autori kao glavnu metriku koriste tačnost, koja je na skupu s
0,172\,\% prevara slabo informativna.

\textbf{Zaključak i značaj za istraživanje.}
Rad potvrđuje da je Random Forest pouzdan izbor za detekciju prevara
i da Naive Bayes ima najlošiji odnos preciznosti i osjetljivosti.
Za naš projekat je važna i pouka o evaluaciji: kod nebalansiranih
podataka rezultate treba prikazivati kroz preciznost, recall i
F1-mjeru za klasu prevare, a ne samo kroz tačnost.


\Needspace{12\baselineskip}
\subsubsection*{Naučnoistraživački rad -- \ClanPet}

\textbf{Naslov rada:}
\emph{Credit Card Fraud Detection Using Improved Deep Learning Models}.\\

\textbf{Autori:}
Sumaya S. Sulaiman, Ibraheem Nadher i Sarab M. Hameed.\\

\textbf{Izvor:}
\emph{Computers, Materials \& Continua}, vol. 78, br. 1, str. 1049--1069,
2024. \cite{sulaiman2024fraud}

\begin{table}[H]
\centering
\caption{Sažeti pregled naučnoistraživačkog rada Sulaiman i saradnika (2024)}
\label{tab:sulaiman-pregled}
\small
\renewcommand{\arraystretch}{1.22}
\begin{tabularx}{\linewidth}{@{}p{0.25\linewidth}Y@{}}
\toprule
\textbf{Stavka} & \textbf{Opis} \\
\midrule

Problem &
Detekcija prevara kreditnim karticama pomoću modela dubokog učenja
i ispitivanje uticaja podešavanja hiperparametara i metoda
uravnotežavanja na njihove performanse. \\

Skup podataka &
Evropski skup transakcija kreditnim karticama (Kaggle, septembar 2013.):
284\,807 transakcija, od kojih su 492 prevarne (0,172\,\%). \\

Korištene metode &
Autoenkoder (AE), konvolucijska neuronska mreža (CNN) i LSTM mreža;
podešavanje hiperparametara metodama \emph{random search} i Bayesovom
optimizacijom (cilj: F1-mjera); uravnotežavanje metodama RUS, SMOTE
i ADASYN; stratifikovana unakrsna validacija. \\

Metrike &
Tačnost, stopa detekcije (\emph{detection rate}, tj. recall) i AUC;
F1-mjera kao funkcija cilja pri podešavanju hiperparametara. \\

Najbolji rezultat &
\textbf{LSTM + SMOTE} (\emph{random search}): tačnost \textbf{99,2\,\%},
stopa detekcije \textbf{93,3\,\%}, AUC \textbf{96,3\,\%}. \\

\bottomrule
\end{tabularx}
\end{table}

\textbf{Cilj i istraživački problem.}
Autori polaze od toga da uspješnost modela dubokog učenja u velikoj
mjeri zavisi od izbora hiperparametara. Cilj je ispitati kako
sistematsko podešavanje (broj neurona, \emph{dropout}, aktivacijska
funkcija, veličina \emph{batch}-a, optimizator, funkcija gubitka i
stopa učenja) u kombinaciji s različitim metodama uravnotežavanja
utiče na detekciju prevara.

\textbf{Skup podataka i priprema.}
Korišten je javni evropski skup sa 30 numeričkih atributa i izrazitom
nebalansiranošću klasa. Podaci su podijeljeni na skup za treniranje,
validaciju i testiranje stratifikovanom unakrsnom validacijom, a
uravnotežavanje (RUS, SMOTE ili ADASYN) primijenjeno je \emph{samo
na skupu za treniranje}. Za svaki model izvedeno je osam pokretanja
(četiri slučajnom pretragom i četiri Bayesovom optimizacijom), uz
četiri scenarija: bez uravnotežavanja, RUS, SMOTE i ADASYN.

\textbf{Poređenje eksperimentalnih rezultata.}
Tabela~\ref{tab:sulaiman-rezultati} prikazuje dio rezultata iz Tabele 3
rada (podešavanje hiperparametara slučajnom pretragom).

\begin{table}[H]
\centering
\caption{Rezultati modela dubokog učenja iz rada Sulaiman i saradnika (u \%)}
\label{tab:sulaiman-rezultati}
\small
\setlength{\tabcolsep}{4pt}
\renewcommand{\arraystretch}{1.18}
\begin{tabularx}{\linewidth}{@{}Y c c c@{}}
\toprule
\textbf{Model -- uravnotežavanje} & \textbf{Acc.} & \textbf{Stopa detekcije} & \textbf{AUC} \\
\midrule
AE -- bez uravnotežavanja & 95,0 & 90,0 & 92,7 \\
AE -- SMOTE & 96,1 & 89,7 & 92,9 \\
CNN -- bez uravnotežavanja & 99,7 & 83,0 & 91,4 \\
CNN -- SMOTE & 99,8 & 89,7 & 94,7 \\
LSTM -- bez uravnotežavanja & 99,9 & 80,8 & 90,4 \\
LSTM -- RUS & 98,4 & 88,9 & 93,6 \\
\textbf{LSTM -- SMOTE} & \textbf{99,2} & \textbf{93,3} & \textbf{96,3} \\
LSTM -- ADASYN & 99,2 & 92,6 & 95,9 \\
\bottomrule
\end{tabularx}
\end{table}

\textbf{Tumačenje rezultata.}
Uravnotežavanje skupa za treniranje značajno povećava stopu detekcije:
kod LSTM mreže ona raste s 80,8\,\% (bez uravnotežavanja) na 93,3\,\%
uz SMOTE, uz mali pad tačnosti (s 99,9\,\% na 99,2\,\%). Metode koje
generišu sintetičke primjere (SMOTE, ADASYN) bile su uspješnije od
RUS-a, koji odbacuje dio podataka većinske klase. Autori LSTM i CNN
modelima pripisuju prednost jer mogu naučiti sekvencijalne obrasce
transakcija. Rad ne prikazuje preciznost, pa nije moguće procijeniti
broj lažnih alarma koji prati visoku stopu detekcije.

\textbf{Zaključak i značaj za istraživanje.}
Rad pokazuje da podešavanje hiperparametara i oversampling skupa
za treniranje mogu značajno povećati osjetljivost modela, što je
direktno povezano sa zadacima 2.2 našeg projekta (\emph{tuning} i
uravnotežavanje klasa). Istovremeno, potvrđuje važnost prikaza
preciznosti uz recall: bez preciznosti visoka stopa detekcije ne
govori koliko je model upotrebljiv u praksi.


\subsubsection*{Zbirni pregled radova}
\begin{table}[H]
\centering
\caption{Zbirni pregled analiziranih naučnoistraživačkih radova}
\label{tab:radovi-zbirno}
\footnotesize
\renewcommand{\arraystretch}{1.2}
\begin{tabularx}{\linewidth}{@{}>{\raggedright\arraybackslash}p{.15\linewidth}>{\raggedright\arraybackslash}p{.14\linewidth}Y>{\raggedright\arraybackslash}p{.16\linewidth}>{\raggedright\arraybackslash}p{.22\linewidth}@{}}
\toprule
\textbf{Rad} & \textbf{Skup podataka} & \textbf{Metode} & \textbf{Metrike} & \textbf{Najbolji rezultat} \\
\midrule
Afriyie i sar.\ (2023) \cite{afriyie2023fraud} &
Sparkov &
LR, DT, RF; RUS &
Acc, Prec, Rec, F1, Spec, AUC &
RF: AUC 0,989, recall 0,97, F1 0,17 \\
Al-Bulushi i sar.\ (2026) \cite{albulushi2026fraud} &
Sparkov + 5 neviđenih skupova &
LR, NB, DT, RF, ANN, XGBoost; RUS, ROS, SMOTE; ansambli &
Acc, Prec, Rec, F1, AUC &
XGBoost: F1 0,84; Bagging~1: F1 0,795, recall 0,864 \\
Breskuvienė i Dzemyda (2024) \cite{breskuviene2024fraud} &
Sparkov, sintetički i evropski skup &
RF, XGBoost, CatBoost; selekcija atributa FID-SOM &
F1, MCC, G-Mean, AUC-PR, AUC-ROC &
XGBoost + FID-SOM (Sparkov): F1 0,855, AUC-PR 0,895 \\
Ileberi i sar.\ (2022) \cite{ileberi2022fraud} &
Evropski skup &
DT, RF, LR, ANN, NB; GA selekcija atributa; SMOTE &
Acc, Prec, Rec, F1, AUC &
GA-RF: Acc 99,98\,\%, F1 82,41\,\% \\
Sulaiman i sar.\ (2024) \cite{sulaiman2024fraud} &
Evropski skup &
AE, CNN, LSTM; tuning hiperparametara; RUS, SMOTE, ADASYN &
Acc, stopa detekcije, AUC &
LSTM + SMOTE: stopa detekcije 93,3\,\%, AUC 96,3\,\% \\
\bottomrule
\end{tabularx}
\end{table}

\subsubsection*{Očekivanja na osnovu literature}
\textbf{Algoritmi.} U svim radovima koji ih porede najbolje rezultate daju ansambli
zasnovani na stablima odlučivanja. Random Forest je najbolji model kod Afriyie i
saradnika \cite{afriyie2023fraud} i Ileberija i saradnika \cite{ileberi2022fraud},
a XGBoost i Random Forest kod Al-Bulushi i saradnika \cite{albulushi2026fraud} i
Breskuvienė i Dzemyde \cite{breskuviene2024fraud}. Ovi modeli dobro rade s
mješovitim tipovima atributa, nelinearnim odnosima i atributima različitih skala.
Neuronske mreže postižu konkurentne rezultate \cite{albulushi2026fraud,sulaiman2024fraud},
ali zahtijevaju pažljivo podešavanje hiperparametara. Logistička regresija i Naive Bayes
dosljedno daju najslabije rezultate: na Sparkov skupu F1-mjera im je ispod 0,1
\cite{albulushi2026fraud}, a Naive Bayes ima preciznost svega 7--16\,\% i na
evropskom skupu \cite{ileberi2022fraud}. Ove modele zato očekujemo kao donju
granicu (\emph{baseline}) za poređenje. Kombinovanje modela u ansambl
(\emph{bagging}, \emph{stacking}) daje stabilne rezultate i visok recall, ali ne
nadmašuje uvijek najbolji pojedinačni model \cite{albulushi2026fraud}.

\textbf{Uravnotežavanje klasa.} Uravnotežavanje skupa za treniranje (RUS, ROS, SMOTE)
povećava osjetljivost modela \cite{sulaiman2024fraud,albulushi2026fraud}, ali na
realističnoj raspodjeli klasa značajno smanjuje preciznost
\cite{afriyie2023fraud,albulushi2026fraud}. Primjer je Random Forest treniran na
\emph{undersampling} uzorku: recall 0,97 uz preciznost 0,09--0,17. Za naš projekat
to znači da modele treniramo i na balansiranom skupu, ali ih provjeravamo i na
stratifikovanim uzorcima s originalnim udjelom prevara (0,58\,\%), te da
uravnotežavanje primjenjujemo isključivo na skupu za treniranje.

\textbf{Metrike.} Ukupna tačnost nije pouzdana mjera jer model koji sve transakcije
proglasi legitimnim postiže tačnost od oko 99,4\,\%. Gotovo svi radovi koriste
preciznost, osjetljivost (recall) i F1-mjeru za klasu prevare, te AUC-ROC. Breskuvienė
i Dzemyda \cite{breskuviene2024fraud} pokazuju da je AUC-ROC na nebalansiranim
podacima visok (0,97--0,998) i za osjetno slabije modele, pa preporučuju AUC-PR i MCC.
Kao glavne metrike projekta zato koristimo \textbf{F1-mjeru, preciznost i recall za
klasu prevare, AUC-PR i MCC}, a dodatno tačnost, specifičnost, kappa statistiku
i AUC-ROC, u skladu sa specifikacijom zadatka 2.2.

\textbf{Atributi.} Na Sparkov skupu uočeni su obrasci prevara povezani s iznosom
transakcije, dobom dana i kategorijom kupovine \cite{afriyie2023fraud}, dok
selekcija atributa i atributi izvedeni iz ponašanja korisnika dodatno poboljšavaju
rezultate \cite{breskuviene2024fraud,albulushi2026fraud}. Očekujemo da će izvedeni
atributi (sat transakcije, starost korisnika, udaljenost korisnika od trgovca) biti
korisniji od identifikatora i ličnih podataka, koje treba ukloniti.

\subsubsection*{Istraživačka hipoteza}
\textbf{Radna hipoteza:} \emph{Na korištenom Sparkov skupu podataka modeli zasnovani na
ansamblima stabala odlučivanja (Random Forest, gradient boosting) postići će veću F1-mjeru
i AUC-PR za klasu prevare od linearnih i probabilističkih modela (logistička regresija,
Naive Bayes). Uravnotežavanje skupa za treniranje i podešavanje hiperparametara povećaće
osjetljivost modela, ali uz pad preciznosti na skupu s realističnim udjelom prevara.
Ansambl koji kombinuje različite klasifikatore postići će F1-mjeru najmanje jednaku
najboljem pojedinačnom modelu, uz stabilnije rezultate kroz unakrsnu validaciju.}

Hipoteza će se provjeriti u zadacima 2.2 i 2.3 poređenjem metrika svih modela
na istom testnom skupu i kroz unakrsnu validaciju.
